% results.tex (국문 v10). 차이 = 앞 방법의 RMSE − 뒤 방법의 RMSE, 괄호 안은 95% 블록 부트스트랩 구간, 음수면 앞 방법이 낫다.
\section{결과}\label{sec:results}

\subsection{관측 자료와 관측 연도 기후 입력}
라벨은 알래스카 \nncyAk{}개 셀-연도(\nncellAk{}개 셀, \nyrsAkko), 레나 델타 \nncyLe{}개(\nncellLe{}개 셀, \nyrsLeko), 캐나다 북서부 \nncyCa{}개(\nncellCa{}개 셀, \nyrsCako)이다(그림~1a,b, 표~1). 관측 연도는 고르지 않다. 알래스카 셀-연도의 \ntopyrshareAk\%가 \ntopyrAk{}년에 관측되었고, 레나 델타 라벨은 \nnyrLe{}개 해에 퍼져 있다. 관측은 대부분 8월, 융해가 끝나기 전에 이루어졌다(보충 그림 S1; 관측일 중앙값 알래스카 \ndoymedAk, 레나 델타 \ndoymedLe, 캐나다 \ndoymedCa{}번째 날). 따라서 관측일까지 누적한 융해 도일은 계절 전체 값보다 작고, 그 비의 중앙값은 \ntddfracAk, \ntddfracLe, \ntddfracCa{}이다(그림~1d). 이 연구의 모든 모델은 2010--2024년 ERA5-Land 월자료에서 관측 연도의 기후를 관측일까지 누적해 입력으로 쓴다. 대조 입력은 같은 변수의 2011--2024년 평균이며, 기존 제품의 다년 입력에 해당한다(그림~1c,d). 물리 잔차 워크플로, 비교한 네 모델 구조, 다섯 검증 설계는 그림~2에 정리하였다.

평년 기후 입력을 관측 연도 기후 입력으로 바꾸면 레나 델타에서 재보정 Stefan 모델(\ndforcLePo{} cm)과 물리 잔차 모델(\ndforcLeRo{} cm)의 공간 블록 오차가, 캐나다에서 물리 잔차 모델의 오차(\ndforcCaRo{} cm)가 줄었다(그림~3a). 알래스카에서는 변화를 판정할 수 없었다(Stefan \ndforcAkPo{} cm). 라벨 대부분이 한 해에 몰려 있어 관측 연도 입력이 공간 분포 이상의 정보를 거의 더하지 못하기 때문이다. 관측 연도 입력은 연도가 명시된 지도의 전제이기도 하다. 모든 예측이 명시된 해의 기후를 가리킨다.

\subsection{검증 설계와 학습에서 뺀 지역으로의 전이}
방법의 순위는 검증 설계에 따라 달랐다(그림~3b). 셀 단위 무작위 교차검증에서 직접 ML의 오차가 가장 낮았고(알래스카 \nrmseAkrandDz{} cm), 0.5$^\circ$ 블록의 공간 블록 교차검증에서는 직접 ML의 오차가 알래스카 \noptimAkDz, 레나 델타 \noptimLeDz, 캐나다 \noptimCaDz{} cm 커졌다. 연도 전체를 빼면 라벨이 여러 해에 퍼진 레나 델타에서 직접 ML의 오차가 \nrmseLeblkDz{} cm에서 \nrmseLeyrDz{} cm로 커졌고, 블록과 연도를 함께 빼는 설계는 세 지역 모두에서 모든 방법의 지역 안 오차가 가장 컸다. 한 지역 전체를 빼고 다른 두 지역으로 학습하면 레나 델타에서 직접 ML의 오차가 \nrmseLeregDz{} cm로 커졌고, 다른 지역의 계수를 쓴 Stefan 모델은 \nrmseLeregPz{} cm였다(차이 \ndrefLeregDz{} cm). 같은 시험에서 물리 잔차 모델은 Stefan 모델 수준을 유지하였고(\ndrefLeregRo{} cm), 캐나다에서는 Stefan 모델보다 나았다(\ndrefCaregRo{} cm). 이하에서는 공간 블록 교차검증을 주 설계로, 지역 홀드아웃을 전이 시험으로 쓴다.

\subsection{물리 잔차 모델의 이득과 조건}
관측 연도 입력과 2010--2024년 라벨로 공간 블록 교차검증을 하면, 물리 잔차 모델은 재보정 Stefan 모델의 오차를 레나 델타에서 \nrmseLeblkPo{} cm에서 \nrmseLeblkRo{} cm로(\ndrefLeblkRo{} cm), 캐나다에서 \nrmseCablkPo{} cm에서 \nrmseCablkRo{} cm로(\ndrefCablkRo{} cm) 줄였다. 알래스카에서는 두 모델이 약 \nrmseAkblkRo{} cm로 같았다(\ndrefAkblkRo{} cm). 이 이득에는 지형과 토양 공변량이 필요하였다(그림~3c). 기후 공변량만 쓰면 잔차 모델은 알래스카(\ndfeatAkFclim{} cm)와 캐나다(\ndfeatCaFclim{} cm)에서 Stefan 모델을 개선하지 못하였고, 지형과 토양을 더해야 위의 이득이 생겼다. MODIS 식생 · 적설 시기(2010--2024년)나 고해상 토양 · 피복 분류를 더한 효과는 레나 델타와 캐나다에서 0.5 cm 미만이었고, 알래스카에서는 고해상 토양 · 피복 분류가 작은 이득을 더하였다(\ndfeatAkFhires{} cm).

물리 정보가 들어가는 위치도 결과를 바꾸었다(보충 그림 S2). Stefan 모델의 유사라벨을 학습 자료에 더하면 Stefan 모델을 재현하는 데 그쳤고(레나 델타 \ndrefLeblkDo{} cm), Stefan 예측을 입력으로 쓰면 직접 ML과 비슷하였다. 지역 안 라벨을 모두 쓰면 레나 델타에서 직접 ML의 점추정 오차가 물리 잔차 모델보다 낮았지만(\nrmseLeblkDz{} cm 대 \nrmseLeblkRo{} cm), 레나 델타와 알래스카에서 Stefan 대비 구간이 0을 포함하였고 지역 홀드아웃에서는 크게 나빠졌다. CatBoost 대신 랜덤 포레스트, 다층 퍼셉트론, TabM, FT-Transformer를 써도 레나 델타와 캐나다에서 물리 잔차 이득의 부호는 바뀌지 않았다.

여덟 공변량의 지역 안 3분위로 나누면, 이득은 레나 델타(\ngcLecells{}개 3분위 가운데 \ngcLesig{}개에서 구간이 0 제외)와 캐나다(\ngcCacells{}개 가운데 \ngcCasig{}개)의 대부분 조건에서 나타났고, Stefan 모델의 오차가 큰 조건일수록 컸다. 3분위 간 Stefan 오차와 이득의 상관은 레나 델타 \ngcLecorr, 캐나다 \ngcCacorr{}였다(그림~3d). 레나 델타의 가장 따뜻한 TDD 3분위에서 이득은 \ndtddLewarm{} cm였다. 알래스카에서는 어느 3분위도 \ngcAkmaxabs{} cm를 넘게 바뀌지 않았다. 따라서 잔차 모델은 토양과 지형 때문에 계수가 지역 값에서 벗어나는 곳에서 Stefan 계수를 보정하는 역할을 한다. 입력 기여도 이와 맞는다. 잔차 모델의 평균 절대 SHAP 값 가운데 지형이 레나 델타에서 \nshapLeResTerrain\%, 토양이 캐나다에서 \nshapCaResSoil\%를 차지하였다(보충 그림 S3).

\subsection{적은 라벨과 새 지역}
라벨 수 사다리는 라벨이 적을 때의 거동을 보여 준다(그림~4a). 한 지역의 2010--2024년 셀-연도에서 학습 라벨 25개를 뽑으면, 물리 잔차 모델은 같은 라벨로 재보정한 Stefan 모델과 같거나 나았고(알래스카 \nllAkRotwfive, 레나 델타 \nllLeRotwfive, 캐나다 \nllCaRotwfive{} cm), 라벨 100개에서는 세 지역 모두에서 이득의 구간이 0을 제외하였다(\nllAkRohundred, \nllLeRohundred, \nllCaRohundred{} cm). 라벨 25개의 직접 ML은 알래스카(\nllAkDztwfive{} cm)와 캐나다(\nllCaDztwfive{} cm)에서 Stefan보다 점추정 오차가 컸다.

전이 실험에서는 다른 두 지역의 라벨과 대상 지역 라벨 $n$개로 학습한 모델로 대상 지역을 예측하였다(그림~4b). 대상 라벨이 없을 때 레나 델타에서 직접 ML은 \ntlrmseLeDszero{} cm, 원천 계수 Stefan 모델은 \ntlrmseLePzzero{} cm였다. 대상 라벨 25개를 더하면 물리 잔차 모델은 \ntlrmseLeRstwfive{} cm, 직접 ML은 \ntlrmseLeDstwfive{} cm로 낮아졌다. 알래스카에서는 물리 잔차 모델이 모든 $n$에서 가장 낮았다(대상 라벨 0개에서 \ntlrmseAkRszero{} cm). 캐나다에서는 원천 지역으로 학습한 직접 ML이 가장 정확하였고(\ntlrmseCaDszero{} cm), 캐나다 라벨 25개로 Stefan 계수를 다시 맞추면 오차가 \ntlrmseCaPzzero{} cm에서 \ntlrmseCaPntwfive{} cm로 커졌고, 이 기준값을 쓰는 물리 잔차 모델도 \ntlrmseCaRszero{} cm에서 \ntlrmseCaRstwfive{} cm로 커졌다. 캐나다 지점들의 융해 깊이 범위가 넓기 때문이다.

같은 절차를 모델 개발에 쓰지 않은 외부 CALM 지역 8곳(\nextcells{}개 셀, \nyrsext{}년 관측)에 적용하였다(그림~4c, 보충 그림 S4). Wei 등, Liu 등, ESA CCI v5 값이 모두 있는 7개 지역 \nftrows{}개 셀-연도에서, 세 연구 지역으로 학습한 물리 잔차 모델은 지역 라벨 없이 \nftAllRszero{} cm였다. 같은 외부 지역의 라벨 5개를 더하면 \nftAllRsfive{} cm, 25개를 더하면 \nftAllRstwfive{} cm로 낮아졌고, 직접 ML은 라벨 25개에서 \nftAllDstwfive{} cm였다. 라벨 5개에서 물리 잔차 모델은 ALT 관측으로 학습하지 않은 ESA CCI v5(\nftAllCcifivezero{} cm)보다 낮았다(차이 \nftdRsCcifivefive{} cm). Wei 등과 Liu 등은 이 지점들을 학습 자료로 썼으며 \nftAllWeizero{} cm와 \nftAllLiuzero{} cm였다. 이득의 대부분은 Stefan 계수의 재보정에서 나왔고, 잔차 모델은 이 지점들에서 추가 이득을 내지 않았다(라벨 5개에서 \nftdRsPnfive{} cm). 서 · 중부 시베리아와 스발바르의 융해 깊이가 연구 지역의 범위를 넘기 때문이다.

\subsection{같은 셀 · 같은 해의 공개 제품 비교}
제품과 모델은 Wei 등, Liu 등, ESA CCI v4가 모두 관측 연도 값을 가진 셀-연도에서 비교하였고, 이 조건으로 비교 기간은 \nyrscc{}년이다(그림~5a). 레나 델타에서 물리 잔차 모델은 \nccLeRoall{} cm로, ESA CCI v5(\nccLeCcifiveall{} cm), Aalto 등(\nccLeAaltoall{} cm), ESA CCI v4(\nccLeCcifourall{} cm), Liu 등(\nccLeLiuall{} cm), Wei 등(\nccLeWeiall{} cm)보다 낮았다. 제품 학습 지점에서 25 km 넘게 떨어진 셀에서 짝지은 차이는 Wei 등 대비 \ndprodLeweiRomask{} cm, ESA CCI v5 대비 \ndprodLeccifiveRomask{} cm였다. 알래스카(\nccAkRoall{} cm 대 \nccAkWeiall{} · \nccAkLiuall{} cm)와 캐나다(\nccCaRoall{} cm 대 \nccCaWeiall{} · \nccCaLiuall{} cm)에서는 모델과 두 ML 제품이 같은 수준이었다.

지역 간 차이는 제품의 학습 자료와 대응한다. 알래스카 라벨 셀은 Wei 등의 학습 지점에서 중앙값 \ndistweiAkmed{} km 떨어져 있고, Liu 등은 같은 알래스카 · 캐나다 현장 자료의 이전 판으로 학습하였다. 레나 델타 셀은 중앙값 \ndistweiLemed{} km 떨어져 있고(보충 그림 S5), 두 제품 모두 델타 북부의 얇은 활동층을 과대 추정한다(북위 72$^\circ$ 이북 평균 편향 \nbiasLeNorthWei{} cm와 \nbiasLeNorthLiu{} cm). 레나 델타 라벨 없이 학습한 모델도 같은 셀-연도에서 이미 Wei 등보다 정확하였다(\ntlvprmseLeRsWeizero{} cm 대 \ntlvprmseLeWeizero{} cm, \ntlvpLeRsWeizero{} cm).

연도별로 보면 물리 잔차 모델은 셀-연도 30개 이상인 레나 델타의 \nywLen{}개 해(\nyrsywLeko) 모두에서 Wei 등과 Stefan보다 오차가 낮았고, 알래스카(\nyrsywAkko)에서는 \nywAkRoWein{}개 해 가운데 Wei 등 대비 \nywAkRoWeiwin{}개 해, Stefan 대비 \nywAkRoPowin{}개 해, 캐나다에서는 \nywCaRoPon{}개 해 가운데 Stefan 대비 \nywCaRoPowin{}개 해에서 낮았다(그림~5b). 여러 해 관측된 알래스카 셀에서 관측과 예측의 연도 간 변화 상관은 물리 잔차 모델 \nrdeltaAkRo, Wei 등 \nrdeltaAkWei, Liu 등 \nrdeltaAkLiu{}였다. 레나 델타에서는 제품이 연도 간 변화를 더 잘 따랐다(\nrdeltaLeWei{}와 \nrdeltaLeLiu{} 대 \nrdeltaLeRo).

잔차 모델은 제품 위에도 얹을 수 있다. 레나 델타 라벨로 Wei 등의 제품을 선형 재보정하고 같은 잔차 학습기를 더하면 재보정 제품 대비 오차가 \npadLeRpPcal{} cm 바뀌었고, Stefan 기준 모델은 이 제품 기준 모델보다 여전히 정확하였다(\npadLeRoRp{} cm, 보충 그림 S6). 알래스카와 캐나다에서는 두 기준의 차이를 판정할 수 없었다.

\subsection{연도가 명시된 1 km 지도}
알래스카와 레나 델타의 모든 셀-연도로 학습한 물리 잔차 모델을 2011--2024년 각 해의 ERA5-Land 기후로 1 km 격자에 적용하였다(그림~6). 격자는 ESA CCI 영구동토 확률 50\% 이상인 셀이다(알래스카 \nncellmapAk{}개, 레나 델타 \nncellmapLe{}개). 2011--2024년 평균 대비 지역 평균 편차는 알래스카에서 \nanminyrAk{}년 \nanminAk{} cm부터 \nanmaxyrAk{}년 \nanmaxAk{} cm까지, 레나 델타에서 \nanminyrLe{}년 \nanminLe{} cm부터 \nanmaxyrLe{}년 \nanmaxLe{} cm까지였고, 공간 분포도 해마다 달랐다(그림~6b). 레나 델타에서 관측 셀의 평균 편차는 셀-연도 50개 이상인 \nanobsnLe{}개 해(\nanobsyrsLe{}년) 가운데 \nanobssameLe{}개 해에서 예측 지역 편차와 부호가 같았다($r$ = \nanobsrLe). 해마다 관측 셀 구성이 바뀌는 알래스카에서는 일치가 약했고($r$ = \nanobsrAk), 예측의 연도 간 범위(\nanpredrangeAk{} cm)가 관측 셀의 범위(\nanobsrangeAk{} cm)보다 작았다(그림~6c). 적용 가능 영역 안의 셀은 알래스카 \naoaAk\%, 레나 델타 \naoaLe\%이고, 90\% 예측 구간의 반폭은 \nhwAk{} cm와 \nhwLe{} cm이다. 적용 가능 영역 밖의 셀은 주로 산지와 영구동토 지대의 남쪽 가장자리이며 지도에 표시하였다(보충 그림 S7). 보충 그림 S8은 알래스카 2014년과 레나 델타 2016년의 우리 지도와 Wei 등 · Liu 등의 지도를 비교하고, 보충 그림 S9는 잔차 모델이 Stefan 모델에 더하는 보정량을 보여 준다.
